\color{unidarkgreen}{\subsection{Высокочастотная блокировка}}
\color{black}

\begin{enumerate}

%%%%%%%%%%%%%%%%%%%%%%%%%%%%%%%%%%%%%%%%%%%%%%%%%% ОБШИЕ СВЕДЕНИЯ %%%%%%%%%%%%%%%%%%%%%%%%%%%%%%%%%%%%%%%%%%%%%%%%%%
\item Общие сведения

\begin{enumerate}

\item
Высокочастотная блокировка (ВЧБ) является быстродействующей защитой линии с абсолютной селективностью и выполняет функцию основной защиты линии от всех видов КЗ.

\item
ВЧБ реализуется в виде полукомплектов, устанавливаемых на концах линии, каждый из которых включает микропроцессорный терминал релейной защиты и комплект оборудования высокочастотной (ВЧ) связи, обеспечивающий формирование и приём ВЧ-сигналов по каналу связи с использованием внешних приёмопередатчиков(ПП).

\item
Принцип действия ВЧБ основан на сравнении направлений мощности по концам защищаемой линии, определяемых органом направления мощности нулевой последовательности (ОНМНП) при однофазных КЗ и направленным реле сопротивления при междуфазных.
Сравнение направлений мощности осуществляется косвенно – посредством обмена ВЧ-сигналами между полукомплектами. 

\item
ВЧБ содержит две основные группы измерительных органов (ИО): блокирующие – действующие на пуск ВЧ-передатчика, и отключающие – действующие на отключение выключателя.


\item
Измерительные органы ВЧБ представлена на рисунке \ref{pict_vchb:IO_VChB}.

\medskip

\begin{figure}[H]
\centering
\includegraphics[width=0.8\textwidth,height=0.8\textheight,keepaspectratio]{\fbpath/01.01 ВЧ/561.0781 ВЧБ/img/IO_VChB.pdf}
\captionof{figure}{Измерительные органы ВЧБ}
\label{pict_vchb:IO_VChB}
\end{figure}

\medskip

\item 
\textls[-15]{Уставочные параметры измерительных органов ВЧБ представлены в таблице \ref{app_set:vchb}.}


\item
\textls[-10]{В начальный момент повреждения срабатывают блокирующие органы, обеспечивая ускоренный запуск ВЧ-передатчика и передачу блокирующего ВЧ-сигнала на противоположный конец линии.
Место КЗ, внутри или вне защищаемой линии, определяется по отсутствию или наличию постоянного ВЧ-сигнала в канале связи.}

\item
Если повреждение находится в пределах защищаемой линии, то после срабатывания отключающих органов, имеющих более грубые уставки, ВЧ-передатчики каждого полукомплекта останавливаются, при этом блокирующий ВЧ-сигнал снимается, и оба полукомплекта действуют на отключение своих выключателей.
Если повреждение находится вне защищаемой линии, «за спиной» одного из полукомплектов, то в этом полукомплекте отключающие органы не срабатывают, его ВЧ-передатчик остается в работе, и блокирующий ВЧ-сигнал продолжает передаваться по каналу связи, предотвращая срабатывание защиты обоих полукомплектов.

\end{enumerate}


%%%%%%%%%%%%%%%%%%%%%%%%%%%%%%%%%%%%%%%%%%%%%%%%%% БЛОКИРУЮЩИЕ ИО %%%%%%%%%%%%%%%%%%%%%%%%%%%%%%%%%%%%%%%%%%%%%%%%%%
\item Блокирующие измерительные органы

\begin{enumerate}

\item
\textls[-30]{Логика ВЧБ предусматривает ненаправленный пуск ВЧ-передатчика (независимо от направления мощности короткого замыкания).
Команда на пуск ВЧ-передатчика (логический сигнал \textit{«Пуск ВЧ»}) формируется при срабатывании блокирующих органов, которые имеют независимые настройки их уставок.
При однофазных КЗ пуск ВЧ-передатчика осуществляется от блокирующего органа по току нулевой последовательности \textit{«ИО 3I0 блок ВЧБ»}.
При междуфазных КЗ пуск выполняется от органов локальной логики БК (логический сигнал \textit{«Пуск ВЧ от БК»}), а именно:}
\begin{itemize} 
\item по приращению тока прямой последовательности \textit{«ИО dI1 блок ВЧБ»};
\item по приращению тока обратной последовательности \textit{«ИО dI2 блок ВЧБ»};
\item \textls[-35]{по току обратной последовательности с торможением от тока прямой последовательности \textit{«ИО I2т пуск ВЧБ»}}.
\end{itemize}

\item
Пуск ВЧ-передатчика от «ИО 3I0 блок ВЧБ» продлевается на время, регулируемое уставкой \textbf{«Тв ИО 3I0 блок ВЧБ»}.
А минимальное время пуска ВЧ-передатчика от органов БК определяется выдержкой времени \textbf{«Тбл повт. ввода РС откл»}. 

\item
\textls[-20]{Пуск ВЧ-передатчика от блокирующих органов снимается при срабатывании отключающих ИО или при отключении выключателей защищаемой линии, если в течение последних 200 мс не было срабатывания ДЗШ.}

\item
В устройстве также предусматривается возможность ручного пуска ВЧ-передатчика через входной сигнал ФСУ \textit{«Ручной пуск ВЧ»}.
Ручной пуск разрешается только при отсутствии срабатывания блокирующих или отключающих органов.

\item
Кроме того, пуск ВЧ-передатчика осуществляется в следующих случаях:
\begin{itemize} 
\item при любых операциях с выключателями (при оперативных переключениях и в циклах АПВ) от входных сигналов ФСУ \textit{«Коммутация В1»} и \textit{«Коммутация В2»} с выдержкой времени на возврат, равной 250 мс.
Это необходимо для исключения излишней работы защиты из-за разновременности включения фаз выключателя; 
\item при срабатывании блокировки при сквозных токах через ошиновку, если введен программный ключ \textbf{«БСТО ВЧБ»};
\item \textls[-10]{при выявлении неисправности в цепях напряжения, если введен программный ключ \textbf{«Пуск ВЧ при БНН»};}
\item от внешней команды пуска ВЧ-передатчика (входной сигнал ФСУ \textit{«Внешний пуск ВЧ»});
\item при обнаружении неисправности канала связи устройством автоматической проверки канала (АПК), если введен программный ключ \textbf{«Блок. ВЧБ от АПК»};
\item при обнаружении неисправности приемопередатчика, если введен программный ключ \textbf{«Блок. ВЧБ от Неиспр. ПП»};
\item при обнаружении неисправности терминала системой самодиагностики, если введен программный ключ \textbf{«Пуск ВЧ при неиспр. терм.»};
\item \textls[-30]{при выводе терминала или функции ВЧБ из работы, если введен программный ключ \textbf{«Пуск ВЧ при выводе»}.}
\end{itemize}

\item
Пуск ВЧ-передатчика снимается при появлении сигнала останова ВЧ-передатчика (логический сигнал \textit{«Запрет пуска ВЧ»}).
Это обеспечивает надежное срабатывание защиты противоположного конца линии при каскадном отключении.
Сигнал \textit{«Запрет пуска ВЧ»} формируется в следующих случаях:
\begin{itemize}
\item при срабатывании внутренних защит и УРОВ (входной сигнал ФСУ \textit{«Останов ВЧ»});
\item при приеме внешнего сигнала \textit{«Останов ВЧ внеш.»}.
Останов ВЧ-передатчика от данного сигнала блокируется в случае срабатывания ДЗШ (входной сигнал ФСУ \textit{«Прием Ср. ДЗШ»}), действующего на отключение выключателя защищаемой линии.
Блокировка продлевается на время, равное 200 мс. 
\end{itemize}

\item
Однако пуск ВЧ-передатчика при неисправности терминала, неисправности канала связи или приемопередатчика, а также при выводе терминала или функции ВЧБ из работы не блокируется от сигнала \textit{«Запрет пуска ВЧ»}.

\item
Функциональная схема управления ВЧПП (ВЧБ) представлена на рисунке \ref{pict_vchb:VChPP}.

\medskip

\begin{figure}[H]
\centering
\includegraphics[width=0.95\textwidth,height=0.95\textheight,keepaspectratio]{\fbpath/01.01 ВЧ/561.0781 ВЧБ/img/VChPP.pdf}
\captionof{figure}{Функциональная схема управления ВЧПП (ВЧБ)}
\label{pict_vchb:VChPP}
\end{figure}

\medskip

\item
Устройство АПК, входящее в состав ВЧПП, используется для автоматической проверки ВЧ-канала связи.
К устройству от ВЧПП подводятся сигналы \textit{«Контакт АПК»} и \textit{«Неиспр. ВЧПП»}.
Наличие сигнала \textit{«Контакт АПК»}, свидетельствует об отсутствии неисправности в канале связи.
При выявлении неисправности канала связи формируется логический сигнал \textit{«Сигн. Неиспр. КС»}.
При ручном пуске ВЧ-передатчика, срабатывании блокирующих или отключающих органов формируется сигнал \textit{«Запрет АПК»}, блокирующий действие АПК.
Действие АПК также может быть выведено оперативно от внешнего сигнала \textit{«Вывод АПК»}.
При обнаружении неисправности приемопередатчика формируется логический сигнал \textit{«Сигн. Неиспр. ПП»}.
\textls[-15]{При введенных положениях программных ключей \textbf{«Блок. ВЧБ от АПК»} и \textbf{«Блок. ВЧБ от Неиспр. ПП»} дополнительно формируется сигнал \textit{«Блок. от ВЧПП»}, который действует на вывод ВЧБ из работы и пуск ВЧ-передатчика.}

\item
Функциональная схема контроля ВЧ-канала связи представлена на рисунке \ref{pict_vchb:Ctrl VCh VChB}.

\medskip

\begin{figure}[H]
\centering
\includegraphics[width=0.9\textwidth,height=0.9\textheight,keepaspectratio]{\fbpath/01.01 ВЧ/561.0781 ВЧБ/img/Ctrl VCh VChB.pdf}
\captionof{figure}{Функциональная схема контроля ВЧ-канала связи}\label{pict_vchb:Ctrl VCh VChB}
\end{figure}

\medskip

\item
Обеспечивается формирование сигнала \textit{«Вызов»} при приеме ВЧ-сигнала, длительность которого превышает 5 с.

\item 
Уставочные параметры функциональной схемы управления ВЧПП (ВЧБ) и схемы контроля ВЧ-канала связи представлены в таблице \ref{app_set:vchb}.

\end{enumerate}


%%%%%%%%%%%%%%%%%%%%%%%%%%%%%%%%%%%%%%%%%%%%%%%%%% ОТКЛЮЧАЮЩИЕ ИО %%%%%%%%%%%%%%%%%%%%%%%%%%%%%%%%%%%%%%%%%%%%%%%%%%
\item Отключающие измерительные органы

\begin{enumerate}

\item
Подготовка цепи отключения осуществляется от отключающих органов.
Отключающие органы реализованы с возможностью независимой настройки их уставок по:
\begin{itemize} 
\item по сопротивлению \textit{«РС откл ВЧБ»};
\item направлению мощности нулевой последовательности \textit{«ОНМНП откл ВЧБ»};
\item току нулевой последовательности \textit{«ИО 3I0 откл ВЧБ»};
\item напряжению нулевой последовательности \textit{«ИО 3U0 откл ВЧБ»}.
\end{itemize}

\item
При однофазных КЗ подготовка цепи отключения осуществляется от ОНМНП, работа которого подтверждается срабатыванием отключающего органа по току нулевой последовательности.
Предусмотрена возможность блокировки отключающего органа по току нулевой последовательности при выявлении броска намагничивающего тока – вводится программным ключом \textbf{«Блок. ИО 3I0 откл при БНТ»}. 

\item
Угол максимальной чувствительности ОНМНП равен 110°.
Порог чувствительности ОНМНП определяется уставками срабатывания отключающих органов по току и напряжению нулевой последовательности. 

\item
ОНМНП может излишне срабатывать в кратковременном неполнофазном режиме, вызванном неодновременностью замыкания фаз выключателя при его включении, что приведёт к излишнему срабатыванию защиты.
Для предотвращения указанного излишнего срабатывания предусмотрена блокировка действия ОНМНП при включении выключателя – вводится программным ключом \textbf{«Вывод ОНМНП при включении»}.
Время вывода регулируется уставкой \textbf{«Т вывода ОНМНП при включении»}.

\item
При междуфазных КЗ подготовка цепи отключения осуществляется от отключающего РС, для которого предусмотрена отдельная логика блокировки при качаниях.
Ввод в работу отключающего РС (логический сигнал \textit{«Ввод РС откл ВЧБ от БК»}) осуществляется от срабатывания блокирующих органов по приращениям токов прямой или обратной последовательностей или пускового органа по току обратной последовательности с торможением от тока прямой последовательности «ИО I2т пуск ВЧБ» на время, задаваемое уставкой \textbf{«Тввода РС откл»}.
Факт формирования разрешающего сигнала \textit{«Ввод РС откл ВЧБ от БК»} подхватывается RS-триггером на время, определяемое уставкой \textbf{«Тбл повт. ввода РС откл»}, в течение которого запрещается повторный ввод отключающего РС.
Блокировка повторного ввода отключающего РС может быть сброшена при отключении выключателя при введенном программном ключе \textbf{«Сброс БК ВЧБ от РПО»}. 

\item
Характеристика срабатывания «ИО I2т пуск ВЧБ», используемого в логике БК, представлена на рисунке \ref{pict_vchb:IO_I2t}.
Порог срабатывания горизонтального участка задается уставкой \textbf{«Iср I2т пуск ВЧБ»}.
Наклонный участок характеристики проходит через точку, равной $ I_{\textmd{ном}} $ , под углом к горизонтальной оси, соответствующим коэффициенту торможения $ k_{\textmd{т}} $.
Коэффициент торможения задается уставкой \textbf{«Кт I2т ВЧБ»}. 

\medskip

\begin{figure}[H]
\centering
\includegraphics[width=0.4\textwidth,height=0.4\textheight,keepaspectratio]{\fbpath/01.01 ВЧ/561.0781 ВЧБ/img/IO_I2t.pdf}
\captionof{figure}{Характеристика срабатывания ИО тока обратной последовательности с торможением от тока прямой последовательности}\label{pict_vchb:IO_I2t}
\end{figure}

\medskip

\item
Отключающее РС имеет направленную характеристику срабатывания, которая задается уставками \textbf{«Rср РС откл ВЧБ»}, \textbf{«Хср РС откл ВЧБ»} и \textbf{«Ф1 РС откл ВЧБ»}.
Углы наклона характеристики срабатывания органа направленности, обеспечивающего направленность РС, задаются уставками \textbf{«Ф2 ОН РС откл ВЧБ»} и \textbf{«Ф3 ОН РС откл ВЧБ»}.
Для отстройки от нагрузки предусмотрен вырез в характеристике срабатывания, который определяется уставками \textbf{«Rнг РС откл ВЧБ»} и \textbf{«Фнг РС откл ВЧБ»}. 
Работа отключающего РС может быть заблокирована при выявлении неисправности в цепях напряжения при введенном программном ключе \textbf{«БНН РС откл ВЧБ»}.

\item
Характеристика срабатывания отключающего РС аналогична, представленной на рисунке \ref{pict_nvchz:RS_otkl_NVChZ}.

%\medskip
%
%\begin{figure}[H]
%\centering
%\includegraphics[width=0.35\textwidth,height=0.35\textheight,keepaspectratio]{\fbpath/01.01 ВЧ/561.0781 ВЧБ/img/RS_otkl_VChB.pdf}
%\captionof{figure}{Характеристика срабатывания отключающего РС ВЧБ}\label{pict_vchb:RS_otkl_VChB}
%\end{figure}
%
%\medskip

\item
При внешних замыканиях (например, при двухфазном замыкании на землю на обходной связи) распределение токов в защищаемой линии может оказаться таковым, что возможно одновременное действие на останов ВЧ-передатчика от ОНМНП на одном конце линии и отключающего РС на другом.
Отсутствие блокирующего ВЧ-сигнала приведёт к излишнему срабатыванию защиты и отключению неповрежденной линии. 

\item
Для исключения ложной работы защиты в таком режиме предусмотрена блокировка отключающего РС при срабатывании отключающих органов по току и напряжению нулевой последовательности.
Вариант действия блокировки выбирается программным переключателем \textbf{«Блок. РС откл от ИО»}, имеющим 4 положения:
\begin{itemize} 
\item \textit{3I0} – отключающее РС блокируется при срабатывании «ИО 3I0 откл ВЧБ»;
\item \textit{3U0} – отключающее РС блокируется при срабатывании «ИО 3U0 откл ВЧБ»;
\item \textls[-20]{\textit{3I0 или 3U0} – отключающее РС блокируется при срабатывании «ИО 3I0 откл ВЧБ» или «ИО 3U0 откл ВЧБ»;}
\item \textit{3I0 или 3U0} – отключающее РС блокируется при одновременном срабатывании «ИО 3I0 откл ВЧБ» и «ИО 3U0 откл ВЧБ».
\end{itemize}

\item
После срабатывания отключающих органов осуществляется действие ВЧБ на отключение выключателя.
Для согласования разновременности срабатывания блокирующих органов полукомплектов защиты на противоположных концах линии предусмотрена выдержка времени \textbf{«Тсогл ВЧБ»} на подготовку цепи отключения после срабатывания отключающих органов.
Выдержка времени на срабатывание защиты \textbf{«Тср ВЧБ»} выполнена с функцией сброса набора времени при приеме блокирующего ВЧ-сигнала и останова набора времени при кратковременных помехах в ВЧ-сигнале.
Допустимая длительность указанных помех составляет 2 мс.
Выдержка времени \textbf{«Тсогл ВЧБ»} также вносит задержку в цепь останова ВЧ-передатчика от отключающих органов.

\item
Действие ВЧБ на отключение блокируется при следующих условиях:
\begin{itemize} 
\item при срабатывании блокировки при сквозных токах через ошиновку, которая вводится программным ключом \textbf{«БСТО ВЧБ»};
\item при выводе терминала или функции ВЧБ из работы;
\item при обнаружении неисправности канала связи устройством автоматической проверки канала (АПК), если введен программный ключ \textbf{«Блок. ВЧБ от АПК»};
\item при обнаружении неисправности приемопередатчика, если введен программный ключ \textbf{«Блок. ВЧБ от Неиспр. ПП»};
\item при переводе действия ВЧБ только на сигнал (от внешнего сигнала \textit{«Перевод на сигнал»}).
\end{itemize}

\item
Функциональная схема ВЧБ представлена на рисунке \ref{pict_vchb:VChB}.

\medskip

\begin{figure}[H]
\centering
\includegraphics[width=0.9\textwidth,height=0.9\textheight,keepaspectratio]{\fbpath/01.01 ВЧ/561.0781 ВЧБ/img/VChB.pdf}
\captionof{figure}{Функциональная схема ВЧБ}
\label{pict_vchb:VChB}
\end{figure}

\medskip

\item 
Уставочные параметры функциональной схемы ВЧБ представлены в таблице \ref{app_set:vchb}.

\end{enumerate}

\end{enumerate}